\documentclass[11pt]{article}
\usepackage[margin=0.95in]{geometry}
\usepackage{amsmath,amssymb,amsthm}
\usepackage[T1]{fontenc}
\usepackage{lmodern}
\usepackage{microtype}
\usepackage[hidelinks]{hyperref}
\newtheorem{theorem}{Theorem}
\newtheorem{lemma}{Lemma}
\newcommand{\Z}{\mathbb Z}
\newcommand{\N}{\mathbb N}
\newcommand{\conv}[2]{\sum_{i=0}^{#2} #1_i #1_{#2-i}}
\setlength{\parindent}{0pt}
\setlength{\parskip}{0.5em}
\title{The only positive prime value occurs at index 2\\[3pt]
\large A disproof of conjecture 00000007662}
\author{C0ldSmi1e}
\date{October 2026}
\begin{document}
\maketitle
\begin{abstract}
For the formal series defined by the equation in the conjecture, we prove
that $C_n(2)$ is a positive prime precisely when $n=2$, with $C_2(2)=2$.
The prime-index set is therefore finite, contradicting the explicit
infinitude assertion. A Lean 4 project verifies the full argument,
including existence and uniqueness of the original formal series and
the polynomial evaluation at $q=2$.
\end{abstract}

\section{The stated equation and its actual coefficients}
Both language versions of conjecture 00000007662 define the coefficients by
\begin{equation}\label{eq:original}
 F(q,t)=\frac{1-qtF(q,qt)}{1-tF(q,t)}
\end{equation}
and assert that infinitely many indices $n$ have $C_n(2)$ prime, together
with upper and lower bounds for the number of such indices up to $x$.
We use the displayed equation literally: its minus sign and the argument
$t$ in the denominator are retained. Primes mean ordinary positive primes.
The index set is $\N=\{0,1,2,\ldots\}$; starting at 1 makes no difference.

Work first in $\Z[Q][[t]]$. More generally, in any commutative ring $R$
with $q\in R$, the constant term of $1-tF$ is 1, so it is invertible.
Thus \eqref{eq:original} is equivalent to
\begin{equation}\label{eq:cross}
 F=1+t\bigl(F^2-qF(qt)\bigr).
\end{equation}
Here the coefficient of $t^n$ in $F(qt)$ is $q^n$ times that in $F(t)$.
Consequently the coefficient sequence is exactly the triangular recursion
\begin{equation}\label{eq:rec}
 a_0(q)=1,\qquad
 a_{n+1}(q)=\sum_{i=0}^{n}a_i(q)a_{n-i}(q)-q^{n+1}a_n(q).
\end{equation}
Every index on the right is smaller than $n+1$. This constructs a unique
sequence and a unique formal series. Conversely, the series with these
coefficients satisfies \eqref{eq:cross} coefficientwise; invertibility of
$1-tF$ then gives \eqref{eq:original}. In $\Z[Q]$, each coefficient is a
polynomial. Applying evaluation $Q\mapsto2$ to the finite sums and products
in \eqref{eq:rec} gives the integer sequence $a_n=C_n(2)$ with
\[
 a_0=1,\qquad a_{n+1}=\sum_{i=0}^{n}a_i a_{n-i}-2^{n+1}a_n.
\]
No analytic convergence or unproved hypothesis about existence is used.

\newpage
\section{A positive transform and uniform growth}
Define integers $b_n$ by
\begin{equation}\label{eq:b}
 b_0=1,\qquad b_{n+1}=2^{n+1}b_n-\sum_{i=0}^{n}b_i b_{n-i}.
\end{equation}
We prove positivity and growth; neither is imposed as an assumption.

\begin{lemma}\label{lem:product}
Fix $n$. Suppose $u_0=1$, $u_k\ge0$ for $k\le n$, and for every $k<n$,
\[
 2^k u_k\le u_{k+1}\le2^{k+1}u_k.
\]
Then $u_i u_j\le u_{i+j}$ whenever $i+j\le n$.
\end{lemma}
\begin{proof}
Induct on $i$. The case $i=0$ is equality. At the induction step,
$j=0$ is also equality. For $j\ge1$ and $i+1+j\le n$, all bounds below
have index less than $n$, and nonnegativity permits multiplication:
\begin{align*}
 u_{i+1}u_j
 &\le2^{i+1}u_i u_j
 \le2^{i+1}u_{i+j}
 \le2^{i+j}u_{i+j}
 \le u_{i+j+1}.
\end{align*}
The middle comparison uses $i+1\le i+j$.
\end{proof}

\begin{lemma}\label{lem:growth}
For every $n\ge0$,
\begin{equation}\label{eq:growth}
 b_n\ge1,\qquad 2^n b_n\le b_{n+1}\le2^{n+1}b_n.
\end{equation}
\end{lemma}
\begin{proof}
Use strong induction on $n$, proving all three assertions together.
If $n=0$, $b_n=1$. If $n=k+1$, the earlier lower bound and
$b_k\ge1$ give $b_n\ge2^k b_k\ge1$. Thus $b_i\ge0$ for all $i\le n$.
All earlier growth bounds are now available, so Lemma~\ref{lem:product}
gives $b_i b_{n-i}\le b_n$. Hence
\[
 0\le\sum_{i=0}^{n}b_i b_{n-i}\le(n+1)b_n\le2^n b_n.
\]
The elementary bound $n+1\le2^n$ follows by induction starting at $n=0$.
Subtract the displayed convolution bounds from $2^{n+1}b_n$ in
\eqref{eq:b} to obtain the two required growth bounds at $n$.
This proves the induction without using the bounds at $n$ prematurely.
\end{proof}

In particular, for $n\ge3$ the lower bound at $n-1$ gives
\begin{equation}\label{eq:large}
 b_n\ge2^{n-1}b_{n-1}\ge4.
\end{equation}

\newpage
\section{Signs, parity, and the exact prime-index set}
\begin{lemma}\label{lem:sign}
For every $n$, $a_n=(-1)^n b_n$. Thus $a_n<0$ for odd $n$, while
$a_n=b_n\ge1$ for even $n$.
\end{lemma}
\begin{proof}
The case $n=0$ is immediate. Assuming the formula at all indices at most
$n$, every product in the convolution acquires the same factor $(-1)^n$.
Equation~\eqref{eq:rec} at $q=2$ therefore yields
\[
 a_{n+1}=(-1)^n\left(\sum_{i=0}^n b_i b_{n-i}-2^{n+1}b_n\right)
 =(-1)^{n+1}b_{n+1}.
\]
The sign assertions follow from Lemma~\ref{lem:growth}.
\end{proof}

\begin{lemma}\label{lem:even}
Every coefficient with a positive even index is even.
\end{lemma}
\begin{proof}
Write the index as $2k+2$. Reflection $i\mapsto2k+1-i$ pairs the terms
of the preceding odd-index convolution with no fixed point, giving
\[
 \sum_{i=0}^{2k+1}a_i a_{2k+1-i}
 =2\sum_{i=0}^{k}a_i a_{2k+1-i}.
\]
In \eqref{eq:rec}, this even sum is reduced by
$2^{2k+2}a_{2k+1}$, also even.
\end{proof}

\begin{theorem}
For the coefficients of \eqref{eq:original}, evaluated at $q=2$,
\[
 C_n(2)\text{ is a positive prime}\quad\Longleftrightarrow\quad n=2.
\]
\end{theorem}
\begin{proof}
The first three values from \eqref{eq:rec} are $a_0=1$, $a_1=-1$,
and $a_2=2$. Odd indices have negative values by Lemma~\ref{lem:sign}.
For even $n\ge4$, Lemma~\ref{lem:even} gives an even value, and
\eqref{eq:large} gives $a_n=b_n\ge4$. Such a value cannot be a positive
prime, because the only even positive prime is 2. This covers every index.
\end{proof}

It follows that the prime-index set is exactly $\{2\}$. This contradicts
the explicit infinitude conjunct, and therefore disproves the conjecture
as a whole. For real $x\ge0$ its counting function is exactly
\[
 A(x)=\#\{n\le x:C_n(2)\text{ is a positive prime}\}
 =\begin{cases}0,&x<2,\\1,&x\ge2.\end{cases}
\]
Thus the asserted lower bound cannot hold under either the eventual
lower-bound or the infinitely-often convention for $\Omega(\log\log x)$.
The disproof relies on the false infinitude assertion and does not require
formalizing either asymptotic convention or refuting the upper bound.

\newpage
\section{Formalization and reproducibility}
The accompanying Lean project pins Lean 4.19.0 and Mathlib v4.19.0.
Its objects are actual formal power series, rather than a recurrence
given the name of the conjectured sequence. In namespace
\texttt{Conjecture7662}, the principal correspondence is:
\begin{center}
\small
\begin{tabular}{@{}p{0.49\linewidth}p{0.45\linewidth}@{}}
\textbf{Lean declaration} & \textbf{Mathematical role}\\[4pt]
\texttt{FunctionalEquation} & Exact quotient equation, with unit inverse\\[3pt]
\texttt{functionalEquation\_coefficients} & Equivalence with the coefficient recursion\\[3pt]
\texttt{exists\_unique\_series} & Existence and uniqueness over a commutative ring\\[3pt]
\texttt{polynomial\_evaluation} & $\Z[Q]$ coefficients specialize to $a_n$\\[3pt]
\texttt{b\_growth}, \texttt{coefficient\_signed} & Growth bounds and alternating sign\\[3pt]
\texttt{even\_index\_coefficient} & Evenness at every positive even index\\[3pt]
\texttt{positivePrime\_iff} & Exact prime classification of $a_n$\\[3pt]
\texttt{formal\_series\_disproof} & Unique existence and failure of infinitude for every actual solution\\[3pt]
\texttt{polynomial\_prime\_iff} & The same classification for $C_n(2)$\\[3pt]
\texttt{primeCounting\_eq} & Exact count for every real cutoff
\end{tabular}
\end{center}
The predicate \texttt{PositivePrime(z)} is
$\exists p\in\N,\ \operatorname{Prime}(p)\land z=(p:\Z)$.
It does not take an absolute value or truncate negative coefficients.
The final disproof has no additional growth, sign, or primality hypotheses.

The standalone auxiliary Python checker compares two exact integer methods:
the triangular recursion and iteration of the original quotient, using
truncated formal-series inversion and substitution $t\mapsto qt$.
It checks all coefficients through degree 32 for
$q\in\{-2,-1,0,1,2,3,4\}$ (231 comparisons) and verifies each
truncated quotient identity.
The beginning at $q=2$ is
\[
 1,\ -1,\ 2,\ -11,\ 150,\ -4474.
\]
These finite checks provide implementation cross-checks; the all-index
mathematical result is the Lean theorem, not an extrapolation of the tests.

From the submitted \texttt{lean/} directory, reproduce the proof with
\begin{verbatim}
lake exe cache get
lake build
lake env lean -DwarningAsError=true Conjecture7662.lean
lake env lean -DwarningAsError=true Check.lean
\end{verbatim}
The package README gives the auxiliary command and describes all evidence.
Verification records include a fresh independent project build, strict
source replay, declared theorem axioms, a complete inventory of compiled
module declarations, pinned dependency identities, and matching report
source/PDF checks. These are local checks; acceptance belongs to the
competition maintainers.

\textbf{Source.} The official bilingual statement is included byte-for-byte
as \texttt{conjecture.md} and is available in the repository at
\href{https://github.com/The-Last-Math-Competition/The-Last-Math-Competition/blob/main/conjectures/00000007662.md}{\texttt{conjectures/00000007662.md}}.
The argument above is derived directly from that displayed equation.
\end{document}
